\documentclass[11pt]{article}
\usepackage{amsmath,amssymb,amsthm}
\usepackage[hidelinks]{hyperref}
\newtheorem{result}{Result}
\newtheorem{conjecture}{Conjecture}
\newtheorem{theorem}{Theorem}
\newcommand{\card}[1]{\textsuperscript{\texttt{[#1]}}}

\title{The period of the simple pendulum to second order in the amplitude}
\author{Note on session S1, 24 November 2026}
\date{}

\begin{document}
\maketitle

\begin{abstract}
The question was the coefficient $c_2$ of the $\theta_0^2$ correction to the
small-oscillation period $T_0 = 2\pi\sqrt{L/g}$ of the simple pendulum,
derived analytically and checked against a numerical evaluation of the exact
period. The session finds $c_2 = 1/16$, the classical value, from the exact
period (a complete elliptic integral) and again from a Poincar\'e--Lindstedt
expansion. An independent computation evaluates the exact period at 13
amplitudes from 0.01 to 3~rad and recovers $c_2$ from the numerical data
alone. The next coefficient is $c_4 = 11/3072$. The second-order formula
underestimates the period at every amplitude tested: its relative error stays
below $0.4\%$ up to 1~rad and reaches $39\%$ at 3~rad.
\end{abstract}

\section*{Setting}
The pendulum is a point mass on a massless rigid rod of length $L$, moving in
a vertical plane under uniform gravity $g$, without friction. The angle
$\theta(t)$ of the rod from the downward vertical, in radians, obeys
\begin{equation}
  \ddot\theta + \omega_0^2 \sin\theta = 0, \qquad
  \omega_0 = \sqrt{g/L}, \qquad T_0 = \frac{2\pi}{\omega_0} = 2\pi\sqrt{L/g}.
  \label{eq:eom}
\end{equation}
The mass drops out of the period. The amplitude $\theta_0 \in (0,\pi)$ is the
largest angle reached, so the pendulum never goes over the top. The exact period $T(\theta_0)$ is expanded as
$T = T_0\,(1 + c_2\,\theta_0^2 + c_4\,\theta_0^4 + \dots)$.
The complete elliptic integral of the first kind is written with the modulus
$k$,
\begin{equation}
  K(k) = \int_0^{\pi/2} \frac{d\phi}{\sqrt{1 - k^2\sin^2\phi}},
\end{equation}
and every series of $K$ in this note is a series in $k$. SciPy and mpmath take
instead the parameter $m = k^2$ as argument; the computation below passes them
$m = \sin^2(\theta_0/2)$. These are the conventions of \texttt{conventions.md}.

\section{The second-order coefficient}

\begin{result}[Second-order coefficient of the period]
For the simple pendulum at amplitude $\theta_0$,
\begin{equation}
  T(\theta_0) = T_0\left(1 + \frac{\theta_0^2}{16} + O(\theta_0^4)\right),
  \qquad T_0 = 2\pi\sqrt{L/g},
  \label{eq:result}
\end{equation}
so $c_2 = 1/16$.\footnote{Derivation D-S1-1 and the independent computation
C-S1-2, each confirmed by a blind review.}\card{CL-S1-1}
\end{result}

This is the classical value; the session reproduces and checks it.

\subsection{Derivation}
Derivation D-S1-1 starts from energy conservation,
$\dot\theta^2 = 2\omega_0^2(\cos\theta - \cos\theta_0)$. The four quarter
swings of one period take equal times, by the symmetries
$\theta \to -\theta$ and $t \to -t$, so
\begin{equation}
  T(\theta_0) = \frac{4}{\omega_0}\int_0^{\theta_0}
  \frac{d\theta}{\sqrt{2(\cos\theta - \cos\theta_0)}} .
\end{equation}
With $k = \sin(\theta_0/2)$ and the substitution $\sin(\theta/2) = k\sin\phi$,
this becomes the exact period
\begin{equation}
  T(\theta_0) = 4\sqrt{\frac{L}{g}}\;K\bigl(\sin(\theta_0/2)\bigr)
  = \frac{2T_0}{\pi}\,K(k).
  \label{eq:exact}
\end{equation}
Expanding the integrand of $K$ by the binomial series and integrating term by
term (Wallis integrals) gives, for $0 \le k < 1$,
\begin{equation}
  K(k) = \frac{\pi}{2}\sum_{n \ge 0}\left[\binom{2n}{n}\frac{1}{4^n}\right]^2 k^{2n}
  = \frac{\pi}{2}\left(1 + \frac{k^2}{4} + \frac{9k^4}{64} + O(k^6)\right).
  \label{eq:Kseries}
\end{equation}
Inserting $k^2 = \sin^2(\theta_0/2) = \theta_0^2/4 - \theta_0^4/48 + O(\theta_0^6)$
and collecting powers of $\theta_0$,
\begin{equation}
  \frac{T(\theta_0)}{T_0} = 1 + \frac{\theta_0^2}{16} + \frac{11\,\theta_0^4}{3072}
  + O(\theta_0^6).
  \label{eq:series}
\end{equation}
The coefficient $c_2 = \tfrac14\cdot\tfrac14$ comes from the first correction
in $K$ alone. The next one,
$c_4 = -\tfrac14\cdot\tfrac{1}{48} + \tfrac{9}{64}\cdot\tfrac{1}{16} = 11/3072$,
also involves the quartic term of $\sin^2(\theta_0/2)$. The expansion converges
near $\theta_0 = 0$, since $K$ is analytic in $k^2$ for $|k^2| < 1$. The sign
of $c_2$ is the expected one: $\sin\theta < \theta$, so the restoring torque is
weaker than in the linear problem and larger swings are slower.

The derivation also obtains $c_2$ without the elliptic integral, by a
Poincar\'e--Lindstedt expansion of \eqref{eq:eom} with
$\sin\theta \approx \theta - \theta^3/6$. Removing the resonant term at order
$\theta_0^2$ fixes the frequency,
$\omega^2 = \omega_0^2\,(1 - \theta_0^2/8 + O(\theta_0^4))$, hence
$T = 2\pi/\omega = T_0\,(1 + \theta_0^2/16 + O(\theta_0^4))$.

A blind review confirmed the derivation. It reproduced the series
\eqref{eq:series} symbolically, evaluated the period integral by quadrature and
by direct integration of the equation of motion, and checked the
Poincar\'e--Lindstedt step.

\subsection{Check against the exact period}
Computation C-S1-2 was set up without access to the derivation. Symbolically,
it rederives the exact period \eqref{eq:exact} and expands it to order
$\theta_0^8$, and it carries a Poincar\'e--Lindstedt expansion to order
$\theta_0^6$. The two routes agree exactly: $c_2 = 1/16$, $c_4 = 11/3072$ and
$c_6 = 173/737280$. Numerically, it evaluates the exact ratio
$T/T_0 = (2/\pi)K(\sin(\theta_0/2))$ at 13 amplitudes from 0.01 to 3~rad with
mpmath (40 digits), with SciPy and with an arithmetic--geometric-mean
evaluation of $K$; the three agree to within $1.3\times10^{-15}$. A direct
numerical integration of \eqref{eq:eom}, which uses no elliptic function,
agrees with them to within $1.3\times10^{-14}$.

The residual
$T/T_0 - 1 - \theta_0^2/16$ scales as $\theta_0^{4.0014}$ on $[0.01, 0.3]$~rad;
at $\theta_0 = 0.01$, the residual divided by $\theta_0^4$ is $0.0035807526$,
against $c_4 = 0.0035807292$, and the gap is $c_6\,\theta_0^2$. A least-squares
fit of the integration data alone, on $[0.05, 0.5]$~rad, gives
$c_2 = 0.0625000033$; the error of $3\times10^{-9}$ comes from truncating the
fit, not from the integration.

A blind review confirmed the computation. It re-ran the script and reproduced
its output exactly, and it recovered $c_2 = 1/16$ by routes the script does not
use: the hypergeometric form $(2/\pi)K = {}_2F_1(\tfrac12,\tfrac12;1;k^2)$, a
direct expansion of the period integral, a 50-digit quadrature and a different
ODE solver. A three-point Richardson extrapolation of the quadrature at 0.001,
0.002 and 0.004~rad gives $1/16$ to 20 digits.

\subsection{Accuracy of the second-order formula}
Table~\ref{tab:accuracy} compares the exact period with the second-order
formula. The formula underestimates the period at every amplitude tested. Its
relative error stays below $4\times10^{-3}$ up to 1~rad ($57^\circ$), and grows
to $6\%$ at 2~rad and $39\%$ at 3~rad. Near $\theta_0 = \pi$ the exact period
diverges logarithmically, $T/T_0 \approx (2/\pi)\ln\bigl(4/\cos(\theta_0/2)\bigr)$,
which no polynomial in $\theta_0$ can follow.

\begin{table}[ht]
\centering
\begin{tabular}{rrccc}
\hline
$\theta_0$ (rad) & $\theta_0$ (deg) & $T/T_0$, exact & $1 + \theta_0^2/16$ & relative error \\
\hline
0.1 & 5.7   & 1.0006254 & 1.000625 & $-3.6\times10^{-7}$ \\
0.5 & 28.6  & 1.0158525 & 1.015625 & $-2.2\times10^{-4}$ \\
1.0 & 57.3  & 1.0663342 & 1.0625   & $-3.6\times10^{-3}$ \\
1.5 & 85.9  & 1.1619687 & 1.140625 & $-1.8\times10^{-2}$ \\
2.0 & 114.6 & 1.3289045 & 1.25     & $-5.9\times10^{-2}$ \\
2.5 & 143.2 & 1.6429824 & 1.390625 & $-0.15$ \\
3.0 & 171.9 & 2.5712339 & 1.5625   & $-0.39$ \\
\hline
\end{tabular}
\caption{Exact period and second-order formula, from computation C-S1-2. The
relative error is $[T_0(1 + \theta_0^2/16) - T]/T$.}
\label{tab:accuracy}
\end{table}

\section{What did not work}
Nothing failed in this session. One error was anticipated and ruled out: SciPy
and mpmath take the parameter $m = k^2$, so passing them the modulus
$k = \sin(\theta_0/2)$ would give a wrong period. The computation passes
$m = \sin^2(\theta_0/2)$, and the direct integration of the equation of
motion, which uses no elliptic function, agrees with the library values to
$10^{-14}$.

\section{Open question and next steps}
The goal of the project is answered: $c_2 = 1/16$, derived by two analytic
routes and checked against the exact period at 13 amplitudes. The evidence is
robust, since a derivation and an independent computation each passed a blind
review. The result stays a working claim until the lead copies it into the
draft. The project lists no next step, so the direction is for the lead to set.
Two continuations need little work. The exact period \eqref{eq:exact} and
$c_4 = 11/3072$ could get claim cards of their own, so that they can enter the
draft as results; both artefacts and both reviews already check them. A short
literature survey could supply a reference for this classical result, since
the workspace has no paper card and this note cites nothing.

\bibliographystyle{plain}
\bibliography{references}
\end{document}
